\section{A conformal cut of prescribed capacity}\label{c:section}

The comparison argument will require a pointwise length density on its inner
edge.  A lower bound for the total length of every crossing does not directly
provide such a density.  The following lemma constructs both the cut and its
conformal parametrization; its leading coefficient at infinity can be
prescribed independently of the lower length bound.

Write
\[
 \mathcal S=\R\times(0,\pi),\qquad Z=e^{v+i\alpha}.
\]
The sides $\alpha=0$ and $\alpha=\pi$ correspond to the positive and negative
real rays, respectively.  In the lemma, a crossing is a rectifiable path with
ordinary endpoints on these two rays and otherwise lying in the upper
half-plane.  Allowing limiting or piecewise smooth crossings gives the same
lower-bound hypothesis by approximation.

\begin{lemma}[Variable conformal cut]\label{c:cut}
Suppose a meridional length element in the upper half-plane is
\[
 \dd\ell=B(v,\alpha)\abs{\dd(v+i\alpha)}.
\]
Besides the boundary puncture $Z=0$, allow finitely many punctures
$a_1,\ldots,a_N\in\R\setminus\{0\}$.  Assume the following.
\begin{enumerate}[label=\textup{(\roman*)}]
\item Away from the punctures and the two ends, $B/\sin\alpha$ has a
positive continuous extension to the closed strip.
\item As $v\to-\infty$, this quotient converges uniformly to a positive
continuous function $b_-(\alpha)$.  As $v\to+\infty$,
\[
 \frac{B(v,\alpha)}{c_+e^{2v}\sin\alpha}\longrightarrow1
 \quad\hbox{uniformly},\qquad c_+>0.
\]
At the sides, these statements concern the continuous quotients.
\item Near each $a_j$, put $t=-\log\abs{Z-a_j}$ and
$\beta=\arg(Z-a_j)$.  For sufficiently large $t$, the same length element
has the form
\[
 B_j(t,\beta)\abs{\dd(-t+i\beta)},\qquad
 c_j\sin\beta\le B_j(t,\beta)\le C_j\sin\beta,
 \qquad 0<c_j\le C_j<\infty.
\]
\end{enumerate}
If every crossing has length at least $L>0$, then, for every $k>0$ and
$0<L'<L$, there is a holomorphic univalent map on $\abs{\xi}>1$, smooth
with nonvanishing derivative on $\abs{\xi}=1$, such that
\[
 f(\overline\xi)=\overline{f(\xi)},\qquad
 f(\xi)=k\xi+O(1)\quad(\xi\to\infty).
\]
Its boundary is a smooth Jordan curve enclosing $0$, its upper semicircle
maps to a cut avoiding all additional punctures, and, in the coordinates
$\xi=e^{\sigma+i\vartheta}$,
\begin{equation}\label{c:density}
 \left.\frac{\dd\ell}{\dd\vartheta}\right|_{\sigma=0}
 \ge \frac{L'}2\sin\vartheta,
 \qquad 0\le\vartheta\le\pi.
\end{equation}
Additional punctures not enclosed by the cut are marked points on the real
sides of the parametrized exterior.  The constants used to construct $f$
may depend on $k$.
\end{lemma}

Only continuity of the original weight is required.  In particular the
statement applies to continuous weights with piecewise smooth seams.  The
smoothness asserted for the cut is in the coordinate $Z$.

\subsection{A smooth lower weight}

\begin{lemma}\label{c:lower-weight}
Under the hypotheses of Lemma~\ref{c:cut}, for every $\ell<L$ there is a
weight
\[
 W(v,\alpha)=b_0(v,\alpha)\sin\alpha\le B(v,\alpha)
\]
whose crossing lengths are at least $\ell$, where $b_0$ is smooth, has even
reflection at both sides, satisfies
$0<m_0\le b_0\le M_0<\infty$, and is independent of $v$ on each sufficiently
far end.  It is smooth across the locations of the additional punctures.
\end{lemma}

\begin{proof}
All losses below can be made arbitrarily small.  We first remove the extra
puncture tails, while retaining the original weight elsewhere.  Take
pairwise disjoint puncture neighborhoods and modify the weight only where
$t>T$.  A positive bounded lower extension exists there: in the main strip,
the puncture comparison implies
\[
 \frac{B}{\sin\alpha}\asymp
 \frac{\abs{Z}^2}{\abs{Z-a_j}^2}.
\]
Indeed $\sin\beta=\operatorname{Im}Z/\abs{Z-a_j}$ and
$\abs{\dd\log(Z-a_j)}=\abs{Z}\abs{Z-a_j}^{-1}
\abs{\dd\log Z}$.  Thus a small positive smooth quotient can replace the
divergent quotient near $a_j$, with the transition wholly in $t>T$.

These modifications cannot decrease the crossing infimum by a fixed
positive amount as $T\to\infty$.  To prove this, suppose that modified
crossings of length at most a fixed $M$ did produce such a loss.  Enlarge
each original side by all the modified tails attached to that side.
Between the last encounter with the initial enlarged side before the first
encounter with the terminal enlarged side, keep the intervening path.
It meets no modified tail in its interior.  If an endpoint is in a tail,
this path must traverse the unchanged collar $T/2<t<T$ to leave its
puncture neighborhood.  On that traversal,
\[
 M\ge c_j\int\sin\beta\,\abs{\dd t}
 \ge \frac{c_jT}{2}\min\sin\beta.
\]
At a point attaining, or arbitrarily approaching, this minimum, attach the
path to the nearer local ray by a constant-$t$ angular interval.  Its
original length is at most
\[
 C_j(1-\abs{\cos\beta})\le C_j\sin^2\beta
 \le C_j\left(\frac{2M}{c_jT}\right)^2.
\]
Both local rays at $a_j$ lie on the same original global side.  Attaching
at the two ends of the retained path therefore produces an original
opposite-side crossing with only $O(T^{-2})$ additional length.  The
finitely many disjoint collars make this estimate uniform.  This
contradicts a fixed loss.  Taking $M>L$ suffices, since longer paths
cannot violate any proposed lower bound below $L$.

At the negative end, uniform convergence permits replacement by a smooth
even profile below $b_-$, with arbitrarily small relative loss.  A transition
can be made within a region where the original quotient is uniformly
close to $b_-$.  At the positive end, fix a small $\delta>0$ and take $V$
large enough that
\[
 (1-\delta)c_+e^{2v}\le B/\sin\alpha
 \le(1+\delta)c_+e^{2v}\qquad(v\ge V-1).
\]
Blend downward to $(1-\delta)c_+e^{2v}$ before $v=V$, and thereafter use
$(1-\delta)c_+e^{2\min(v,V)}$.  Radial projection
$(v,\alpha)\mapsto(\min(v,V),\alpha)$ converts a path for this lower
weight into a path for the preexisting weight with length increased by
at most $(1+\delta)/(1-\delta)$.  Below $V$ this follows from the same
ratio comparison on the transition; above $V$, projection deletes the
radial component, and the angular weight at $V$ is at most that ratio
times the saturated weight.  Projection preserves the two sides.
Consequently saturation loses only this arbitrarily small factor in the
crossing lower bound.

The resulting positive continuous quotient is bounded above and below,
and is constant in $v$ on both far ends.  Approximate it uniformly from
below by a smooth quotient, even at the sides and still product on the
ends.  Such approximation follows by even reflection, convolution on
the compact remaining rectangle, and a small downward adjustment; the
positive lower bound converts uniform error to relative error.  The
product profiles can be approximated in the same way before extending
them along the tails.  The accumulated relative and additive losses
can be chosen smaller than $L-\ell$.  This proves the lemma.
\end{proof}

\subsection{Distance calibration and removal of critical levels}

\begin{lemma}[Angular calibration]\label{c:angular}
For every $L'<L$ there are a lower weight $W$ as in
Lemma~\ref{c:lower-weight}, a number $L_1>L'$, and a smooth angular
coordinate $\Theta:\overline{\mathcal S}\to[0,\pi]$ such that
\begin{equation}\label{c:calibration}
 \varphi=\cos\Theta,\qquad
 \abs{\dd\varphi}_{\mathrm{flat}}
 \le\frac{2}{L_1}W\le\frac{2}{L_1}B.
\end{equation}
Here $\Theta=0,\pi$ on the respective sides, $\dd\Theta$ never vanishes,
$\Theta$ has ordinary odd reflection at both sides, and on each far end
it is an increasing angular diffeomorphism independent of $v$.
Moreover, $\Theta$ is joined to $\alpha$ by a smooth family of angular
coordinates with these properties, product on uniform far ends; the
estimate in Equation~\eqref{c:calibration} is needed only at the final member.
\end{lemma}

\begin{proof}
Choose a lower weight with crossing infimum at least $\ell>L'$, leaving
strict room for the losses that follow.  Let $d_+$ denote its distance to
the side $\alpha=0$.  The distance inequality gives
$\abs{\dd d_+}\le W$ almost everywhere.  Since
$m_0\le b_0\le M_0$, angular intervals give
\[
 d_+(v,\alpha)\le M_0(1-\cos\alpha),\qquad
 d_+(v,\alpha)\ge\ell-M_0(1+\cos\alpha).
\]
The second inequality follows by appending an angular interval from the
point to the opposite side to an almost minimizing path for $d_+$.
For a small $h>0$, truncate $d_+$ to $[h,\ell-h]$ and put
\[
 \varphi_0
 =1-2\min\left\{1,\max\left\{0,
                  \frac{d_+-h}{\ell-2h}\right\}\right\}.
\]
Thus $\abs{\dd\varphi_0}\le2W/(\ell-2h)$, and $\varphi_0$ is identically
$1$ and $-1$ on full, uniform collars of the respective sides.  Take $h$
small enough to retain strict slack above $L'$.

Clamp $v$ to a large closed interval with endpoints inside the product
tails.  This projection is length nonincreasing for $W$, so composition
preserves the gradient bound and makes $\varphi_0$ independent of $v$
on the far tails.  Smooth after even reflection at the sides.  The
constant side collars remain constant; elsewhere $W$ has a positive
lower bound on a compact rectangle, so convolution and a small slack
loss preserve the weighted gradient bound.  This produces a smooth
function.  A sufficiently small convex combination with $\cos\alpha$
then gives strict interior values and
\[
 \partial_{\alpha\alpha}\varphi(v,0)<0,
 \qquad \partial_{\alpha\alpha}\varphi(v,\pi)>0.
\]
Its gradient cost fits within the slack, because
$\sin\alpha\le W/m_0$.

On a product end write $W=b_\pm(\alpha)\sin\alpha$ and set
\[
 I_\pm=\int_0^\pi b_\pm(s)\sin s\,\dd s,
 \qquad
 \varphi_\pm(\alpha)
 =1-\frac{2}{I_\pm}\int_0^\alpha b_\pm(s)\sin s\,\dd s.
\]
The vertical crossing on that end gives $I_\pm\ge\ell$.
Both the smoothed function and $\varphi_\pm$ have the same endpoint
values and even reflection, so their difference is $O(\sin^2\alpha)$,
uniformly on the end.  Interpolate to $\varphi_\pm$ sufficiently slowly
in $v$.  The angular derivative is bounded by the convex combination
of the preceding bounds.  The new $v$ derivative is bounded by
$C\abs{\chi'}\sin^2\alpha$, which is arbitrarily small relative to $W$
when the interpolation interval is long.  Thus, for some $L_0>L'$,
\begin{equation}\label{c:initial-calibration}
 \abs{\dd\varphi}\le\frac{2}{L_0}W.
\end{equation}
The angular derivative is strictly negative on the far ends and in
small uniform side collars.  A compactly supported arbitrarily small
Morse perturbation away from those collars leaves these properties
intact and leaves only finitely many critical points.  We retain the
notation $\varphi$ and a slightly smaller $L_0>L'$.

We next replace its level sets by a foliation without losing the
estimate.  Delete small open intervals about its finitely many critical
values.  In a closed remaining regular value band, intersect a level
with $[-V,V]\times[0,\pi]$, taking $V$ in the common product tails.
There is exactly one level point on each vertical edge and no point
on a horizontal edge.  A compact regular level is a union of circles
and intervals; its two boundary points must therefore be joined by
one interval component.  This is the unique essential component.

For an entire closed regular band, a vector field satisfying
$\dd\varphi(V_\varphi)=1$, chosen tangent to the vertical edges,
flows these essential components smoothly into one embedded rectangle.
It has precisely the two extreme essential arcs and the appropriate
vertical-edge intervals as its boundary.  The Jordan separation theorem
therefore identifies its image with the full region between the two
arcs.  In particular, discarding the other, circular level components
creates no holes in this essential band.  Essential bands at distinct
values are disjoint and ordered, since their order is fixed at both
vertical ends and disjoint crossing arcs cannot exchange that order.
Near $1$ and $-1$ keep the monotone side collars: outside smaller such
collars the values stay uniformly away from the extrema, by compactness
and the product-end behavior.

These finitely many essential bands can be completed to a foliation of
the whole strip by proper arcs.  Here is the extension with its boundary
conditions.  Slightly shrink each good band, retaining its level
coordinate on a neighborhood of the resulting closed band.  Between
two consecutive bands, their smooth boundary arcs and the two vertical
edges bound a disk.  Prescribe product parametrizations on collars of
all four edges, using the original label where it is already fixed and
the existing angular labels on the vertical ends.  They agree at the
corners because the ends are product.  Shrink the collars so that
nonadjacent ones are disjoint and round within the compatible corner
charts.  Smooth Schoenflies charts reduce the extension to a prescribed
orientation-preserving collar map of a disk.  Here the relative
extension can be described explicitly.  In polar coordinates its
boundary derivative is triangular, with positive diagonal entries.
Interpolation to the radial product map is therefore an isotopy on a
sufficiently thin collar; extend its generating vector field with a
cutoff over the disk.  For the resulting boundary map, choose an
increasing lift $h$ with $h(\theta+2\pi)=h(\theta)+2\pi$.  The map
\[
 (r,\theta)\longmapsto
 \bigl(r,\theta+\chi(r)(h(\theta)-\theta)\bigr),
\]
where $0\le\chi\le1$, $\chi=1$ near the boundary, and $\chi=0$ near the center,
is a smooth disk diffeomorphism: its angular derivative is
$1-\chi+\chi h'>0$, and it is the identity near the center.
Undoing the extended collar isotopy gives an extension agreeing with
the entire prescribed smaller collar.  Thus there is no additional
corner or orientation condition to impose.  Extending by the product
ends gives a smooth value function $\widehat\varphi$, with no interior
critical points, agreeing with $\varphi$ on the smaller good bands and
on the side collars.  It has the same endpoint values and reflection parity.
The labels in every intervening strip lie solely in its assigned
complementary value interval.  Thus discarded circles do not reappear
as extra components of a good label.

There is no a priori gradient estimate in an intervening strip, but
there are only finitely many of them, they avoid the side collars, and
they are product at infinity.  Hence
$\abs{\dd\widehat\varphi}/W$ has a finite bound there, say $M$.
The total width $\delta$ of their label intervals, including small
transition intervals within good bands, can be chosen arbitrarily
small.  Choose a smooth positive function $r$ of the label, equal to
$1$ away from those intervals, at most $1$ everywhere, and at most
$\varepsilon$ on every bad label interval; arrange its transitions
inside good bands.  Put
\[
 F(t)=-1+\frac{2}{\int_{-1}^1r(s)\,\dd s}
                    \int_{-1}^t r(s)\,\dd s.
\]
The normalization factor is at most $2/(2-\delta)$.
On good bands the old gradient estimate is multiplied by at most this
factor.  On bad bands it is bounded by this factor times $\varepsilon M$.
First choose $\delta$ small, and then $\varepsilon$ small.  The function
$F\circ\widehat\varphi$ satisfies Equation~\eqref{c:calibration} for some
$L_1>L'$.  It has no interior critical points.  Since $F'>0$ at the
endpoints, its side jets remain nondegenerate and even.  We again call
this function $\varphi$.

Set $\Theta=\arccos\varphi$.  In the interior this is a submersion.
At $\alpha=0$, the even expansion has the form
$1-\varphi=\alpha^2 a(v,\alpha^2)$ with $a(v,0)>0$; smooth square-root
division shows that $\Theta$ is $\alpha$ times a positive smooth even
function.  The same argument applies to $\pi-\Theta$ at the other side.
Thus $\Theta$ has precisely the stated reflection and nonvanishing
boundary derivative.

For use in the analytic argument we also record an isotopy.  Integrate
$D\Theta/\abs{D\Theta}^2$ from the lower edge, with $\Theta$ as time.
The vector field is vertical on the product ends and has bounded
coefficients in the remaining compact region.  Its flow gives a strip
diffeomorphism $D(v_0,\theta)$ with $\Theta(D(v_0,\theta))=\theta$.
Reflection reverses the flow parameter and gives ordinary even/odd
parity for the two components of $D$.  In particular $D$ is product on
uniform far ends.

Such a diffeomorphism is smoothly isotopic to the identity within this
class.  First straighten its two end angular diffeomorphisms by convex
angular interpolation, whose derivative stays positive, and extend
these collar motions by their generating vector fields.  Next
straighten the parametrizations of the horizontal sides, if necessary,
by convex interpolation of their increasing longitudinal
parametrizations; these motions are compactly supported in $v$.
After this, the derivative on each side is diagonal with positive
entries by reflection.  Linear interpolation to the identity in a
sufficiently thin collar is consequently a local collar isotopy.
Extend its generating field with a symmetric cutoff.  The remaining
map fixes a neighborhood of the entire boundary of a compact disk.
After choosing a smaller fixed boundary collar, the smooth disk
isotopy theorem~\cite[Theorem~4]{Smale1959} joins it to the identity
through orientation-preserving diffeomorphisms fixing that collar.
Reattaching the fixed collars and ends, and taking inverse coordinate
maps, gives the required family of angular submersions.  All these
uses of disk topology concern smooth parametrizations of planar disks;
no metric estimate is required during the isotopy.
\end{proof}

\subsection{The nonlinear analytic boundary problem}

We prove the analytic step explicitly because both its normalization at
infinity and its compactness are needed for arbitrary prescribed $k$.
Extend every angular coordinate by odd reflection and by
$\Theta(v,\alpha+2\pi)=\Theta(v,\alpha)+2\pi$.

\begin{lemma}\label{c:analytic}
Let $\Theta$ be an angular coordinate with the product-end, reflection,
and isotopy properties of Lemma~\ref{c:angular}.  For every $k>0$ there
is a real-symmetric holomorphic function $U$ on $\abs{\xi}>1$, smooth
up to the circle and with $U(\infty)=0$, such that
\begin{equation}\label{c:boundary-equation}
 \Theta\bigl(\log k+\Re U(e^{i\vartheta}),
                   \vartheta+\Im U(e^{i\vartheta})\bigr)=\vartheta
 \qquad(\vartheta\in\R).
\end{equation}
The map $f(\xi)=k\xi e^{U(\xi)}$ is univalent on the exterior, and its
boundary is a smooth regular Jordan curve enclosing zero.
\end{lemma}

\begin{proof}
Let $\Theta_s$, $0\le s\le1$, be the isotopy, with $\Theta_0(v,\alpha)=\alpha$
and $\Theta_1=\Theta$.  For fixed $(s,\vartheta)$, the allowed values
$u=x+iy$ form the line
\[
 \Gamma_{s,\vartheta}
 =\{x+iy:\Theta_s(\log k+x,\vartheta+y)=\vartheta\}.
\]
It is a smooth proper embedded line with horizontal ends.  These lines
vary smoothly, periodically in $\vartheta$, with
$\Gamma_{s,-\vartheta}=\overline{\Gamma_{s,\vartheta}}$.
At $\vartheta=0,\pi$ the line is the real axis.  Their heights and the
locations at which they become horizontal have uniform bounds for
$s\in[0,1]$ and $\vartheta$ modulo $2\pi$.

Orient each line from its left end to its right end.  Its tangent phase
has a smooth real lift $\beta$, obtained by transporting the initial
phase zero through the strip diffeomorphisms of the isotopy.  It is
periodic in $\vartheta$ and odd under reflection.  To check the
potential integer ambiguity, at the left horizontal end its value is
zero.  At the right horizontal end its value is an integer multiple of
$2\pi$; that integer is continuous in the isotopy parameter and is zero
for $s=0$, so remains zero.  The same uniqueness of the lift proves
periodicity and reflection, with value zero on the straight boundary
leaves.  On compact sets of positions, the line data, the phase, and
all their derivatives are uniformly bounded.

\paragraph{Openness.}
Work with real-symmetric exterior analytic traces in $C^{m,\lambda}$,
$m\ge2$, $0<\lambda<1$, normalized by $U(\infty)=0$; the boundary equation
takes values in the odd real $C^{m,\lambda}$ functions.  At a solution,
its linearization in a variation $\dot U$ is
\[
 \Theta_{s,v}\Re\dot U+\Theta_{s,\alpha}\Im\dot U
   =c(\vartheta)\Im(e^{-i\beta(\vartheta)}\dot U),
 \qquad c(\vartheta)>0.
\]
Here $c$ is the magnitude of the label gradient, with sign fixed by the
orientation of the strip.  At a symmetric trace it is even, while
$\beta$ is odd.

The normalized exterior Schwarz problem gives a real-symmetric analytic
function $a$, with $a(\infty)=0$, whose imaginary boundary value is
$\beta$.  The zero-mean compatibility condition is automatic for odd
data.  The multiplier $A=e^a$ is nonvanishing, real-symmetric, and
$A(\infty)=1$.  Substituting $\dot U=AH$ reduces a prescribed
linearized value $g$ to
\[
 \Im H=\frac{g}{c\abs A},\qquad H(\infty)=0.
\]
The right side is again odd.  The Schwarz problem solves it uniquely:
the harmonic conjugate is determined by the periodic Hilbert transform,
and the real constant is fixed at infinity.  It gives a bounded inverse
on the stated H\"older spaces.  Multiplication by $A$ preserves the
normalization.  Smooth composition on the circle now permits the
implicit function theorem.  Since $U=0$ solves the equation at $s=0$,
the set of solvable parameters is nonempty and open.

\paragraph{Uniform range bounds.}
The line family bounds $\Im U$ on the circle, and hence throughout the
exterior by the harmonic maximum principle.  Write
$h_{s,+}(\vartheta)$ for the height of the right horizontal end of
$\Gamma_{s,\vartheta}$.  It is a uniformly smooth odd periodic function.
Let $H_{s,+}$ be its normalized exterior Schwarz solution and set
$V=U-H_{s,+}$.  The functions $H_{s,+}$ are uniformly bounded in every
fixed smooth norm.  If $\Re V$ had an excessively large positive
maximum, it would occur at a point of the circle on an open arc where
$U$ is in the right horizontal tail.  On that arc $\Im V=0$.
The Cauchy--Riemann equations then give zero normal derivative of
$\Re V$ at the maximum, contradicting the boundary Hopf lemma unless
$V$ is constant.  The constant case has $V(\infty)=0$ and cannot give
such a maximum.  One may apply this argument after inversion to the
closed disk, where infinity is an interior point.  Applying the same
argument to the left horizontal ends gives a lower bound for $\Re U$.
Thus every solution, at every $s$, satisfies
\begin{equation}\label{c:range}
 \norm U_{L^\infty(\abs\xi\ge1)}\le C.
\end{equation}
The constant may depend on $k$; the only effect of $\log k$ is to
translate the horizontal thresholds for the line family.

\paragraph{Boundary derivative compactness.}
Suppose instead that a sequence of solutions has
$P_j=\max_{\abs\xi=1}\abs{\partial_\vartheta U_j}\to\infty$.
Choose maximum points $e^{i\vartheta_j}$ and use the half-plane
coordinates
\[
 u_j(z)=U_j\bigl(e^{i\vartheta_j+z/P_j}\bigr),
 \qquad \Re z\ge0.
\]
The analytic function $\xi U_j'(\xi)$ has its maximum modulus on the
circle, so Equation~\eqref{c:range} and the definition of $P_j$ give
\[
 \abs{u_j}\le C,\qquad \abs{u_j'}\le1,
 \qquad \abs{\partial_yu_j(0)}=1.
\]
On the imaginary axis its trace lies in
$\Gamma_{s_j,\vartheta_j+y/P_j}$.  Differentiate the defining equation.
With $\beta_j(y)$ the tangent phase along this trace, the result is
\begin{equation}\label{c:scaled-boundary}
 \Im\bigl(e^{-i\beta_j(y)}\partial_yu_j(iy)\bigr)
       =P_j^{-1}g_j(y).
\end{equation}
For example, the forcing is
$(1-\Theta_{s_j,\alpha})/\abs{D\Theta_{s_j}}$, evaluated at the trace.
The range lies in a fixed compact set, the label gradient is bounded
away from zero there, and $u_j$ is uniformly Lipschitz.  Hence both
$\beta_j$ and $g_j$ have uniform Lipschitz bounds on every fixed boundary
interval.

We justify convergence of derivatives all the way to this boundary.
Extend $\beta_j$ from a slightly larger interval by a fixed cutoff, and
solve the half-plane Schwarz problem with that imaginary trace for an
analytic logarithm.  Its exponential $A_j$ and its inverse have uniform
local $C^\lambda$ bounds for every $0<\lambda<1$: the Schwarz integral
and Hilbert transform send compactly supported uniformly Lipschitz data
to uniformly $C^\lambda$ functions.  The analytic derivative
$D_j=i u_j'$ divided by $A_j$ is bounded and, by
Equation~\eqref{c:scaled-boundary}, has uniformly $C^\lambda$ imaginary trace
on the smaller interval.  Subtract a Schwarz solution for a localized
version of that trace.  The remainder has zero imaginary trace and
extends across the interval by Schwarz reflection.  It is uniformly
bounded on a smaller disk, so its derivatives there obey the interior
Cauchy estimates.  Adding back the Schwarz solution proves uniform
local $C^\lambda$ bounds for $D_j$ up to the straight boundary.
Consequently a subsequence converges locally in $C^1$ on the closed
half-plane to a bounded holomorphic $u$, and
$\abs{\partial_yu(0)}=1$.

After a further subsequence $s_j\to s_\infty$ and
$\vartheta_j\to\vartheta_\infty$ modulo $2\pi$.  The limiting boundary
trace lies on the single proper embedded line
$\Gamma_{s_\infty,\vartheta_\infty}$.  Properness and the bound
$\abs u\le C$ place this trace in one compact simple subarc $A$ of that
line: take a compact parameter interval containing the inverse image
of its intersection with the closed range disk.  For every polynomial
$p$, the bounded half-plane maximum principle gives
\[
 \abs{p(u(z))}\le\max_{a\in A}\abs{p(a)}.
\]
There is no boundary condition missing at infinity here.  Map the
half-plane to a disk; the possible exceptional boundary point has
harmonic measure zero, and its contribution disappears by boundedness
when a surrounding boundary arc shrinks.  The displayed inequalities
put the image of $u$ in the polynomial hull of $A$.  In one complex
variable the polynomial hull of a compact set is the set together with
its bounded complementary components.  A planar simple arc has
connected complement, so its hull is itself.  The open mapping theorem
therefore forces $u$ to be constant, contradicting its normalized
boundary derivative.  The assumed derivative blowup is impossible.

\paragraph{Closedness.}
With first derivatives bounded, the same phase-multiplier argument on
fixed boundary intervals gives uniform $C^\lambda$ bounds for them.
The phase and forcing, now smooth compositions of $C^{1,\lambda}$
traces, have $C^{1,\lambda}$ bounds.  More explicitly, on the fixed
circle set $D=i\xi U'$ and let $A$ be the phase multiplier.  Then
$\Im(D/A)=g/\abs A$, with
$g=(1-\Theta_{s,\alpha})/\abs{D\Theta_s}$ evaluated at the trace.
Here $D/A$ is anti-real-symmetric and $g$ is even; this bootstrap
uses Schwarz estimates for unrestricted normalized analytic traces,
not the odd codomain used for openness.  The imaginary boundary mean
vanishes automatically because this actual analytic function has
$(D/A)(\infty)=0$, and its remaining real constant is zero by the same
normalization.  Thus the Schwarz estimates increase the regularity of
$U$ by one derivative.  Induction gives bounds of every fixed order;
interior bounds follow from Cauchy estimates.  Thus sequences of solutions have convergent subsequences
in the chosen $C^{m,\lambda}$ space, and their limits solve the boundary
equation at the limiting parameter.  Solvability is closed as well as
open.  Connectedness of $[0,1]$ proves existence at $s=1$.

\paragraph{Univalence.}
For $f=k\xi e^U$, differentiating the boundary equation gives
\[
 D\Theta\cdot
 \bigl(\partial_\vartheta\Re U,
                  1+\partial_\vartheta\Im U\bigr)=1.
\]
Thus its boundary derivative never vanishes.  If two boundary values
of $f$ coincide, their logarithmic radii agree and their lifted angles
differ by an integer multiple of $2\pi$.  The equivariance of $\Theta$
then says that their parameters differ by the same multiple of $2\pi$.
The boundary curve is therefore injective.  Its winding about zero
is one, directly from $f=k e^{i\vartheta}e^{U(e^{i\vartheta})}$.
It is a regular Jordan curve enclosing zero.

For a point $w$ off this curve, apply the argument principle to the
exterior domain, including its one simple pole at infinity.  If
$n(w)$ is the winding of the positively parametrized boundary curve,
the number of finite preimages, counted with multiplicity, is
$1-n(w)$.  It is one in the unbounded component and zero in the bounded
component.  No interior point can map to the boundary curve, by the
open mapping theorem and the zero count on its bounded side.  Hence
$f$ is univalent, has no interior critical points, and maps onto the
unbounded component.  Symmetry implies that a real image point has a
real preimage, by uniqueness of preimages.  Since the upper half maps
to the upper half near infinity, it does so throughout.  Finally
$U(\infty)=0$ gives $f=k\xi+O(1)$ with precisely the prescribed leading
coefficient.
\end{proof}

\subsection{Completion of the cut construction}

\begin{proof}[Proof of Lemma~\ref{c:cut}]
Choose the angular calibration with $L_1>L'$ and solve the analytic
boundary problem at the prescribed $k$.  On the upper inner edge,
\[
 \varphi\bigl(\log\abs{f(e^{i\vartheta})},
                    \arg f(e^{i\vartheta})\bigr)=\cos\vartheta.
\]
Consequently Equation~\eqref{c:calibration} gives
\[
 \sin\vartheta
 \le \frac{2}{L_1}
 W\bigl(\log\abs f,\arg f\bigr)
 \abs{\frac{\dd}{\dd\vartheta}\log f(e^{i\vartheta})}
 \le \frac{2}{L_1}\frac{\dd\ell}{\dd\vartheta}.
\]
This already proves Equation~\eqref{c:density} with strict slack.  If an endpoint
coincides with an additional puncture, translate $f$ by a sufficiently
small real constant.  This preserves symmetry, univalence, and the
leading coefficient $k$, and preserves enclosure of zero because the
untranslated Jordan curve has positive distance from zero.

The slack survives such a translation, including at the endpoints.
Indeed, in $Z$ coordinates the smooth lower length density is
\[
 \lambda_0(Z)
 =\frac{b_0(\log\abs Z,\arg Z)\operatorname{Im}Z}{\abs Z^2}.
\]
Away from zero it is smooth with odd reflection and has a positive
simple zero at the real axis.  For the translated smooth symmetric
curve, the quotient
\[
 \frac{\lambda_0(f(e^{i\vartheta})+t)
                  \abs{\partial_\vartheta f(e^{i\vartheta})}}
      {\sin\vartheta}
\]
extends continuously to $\vartheta=0,\pi$ and depends continuously on
the small real translation $t$.  At $t=0$ it is at least $L_1/2$.
Compactness of $[0,\pi]$ therefore preserves the bound $L'/2$ for all
sufficiently small $t$.  Avoiding the finitely many translations that
put either endpoint at an $a_j$ makes the cut miss every extra
puncture.  Its interior already lies strictly in the upper half-plane.
The original density dominates this lower density throughout, so the
required estimate follows for $B$.  The statement about remaining
marked side points follows from symmetry and exterior univalence.
\end{proof}
